\chapter{Introduction}
\label{ch:gaussian-introduction}

This supplement asks what the single-point theory of the main manuscript
becomes when every belief and model law is multivariate Gaussian. It cites the
main manuscript and adopts its standing conventions
(\Cref{sec:intro-conventions}), and it states no result that the main
manuscript uses. The belief fiber has dimension \(K\) and the model fiber
dimension \(d_m\), which are not assumed equal. The represented frame groups
are \(\Gamma_b=\rho_b(G)\) and \(\Gamma_m=\rho_m(G)\), with the specialization
\(\Gamma_b=\GL^+(K)\) named wherever it is used
(\Cref{hyp:gen-lg-frame-group}); both channels of an agent use one common
frame, the gauge choice of the main manuscript. The belief-sector precision
studied throughout is that of \(p(k\given m,o,X)\), the fast belief at frozen
model latents (\Cref{sec:gen-lg-frame-group}); its membership in the
interaction family requires \(A_i+O_i\succeq0\) at every agent, so the balance
condition on the prior alone is sufficient but not necessary
(\Cref{prop:gauss-global-interaction-condition}).

The Gaussian supplies explicit determinant gaps, Fisher charts,
matrix-weighted interaction operators, aggregation formulas, Kron
counterexamples, and generalized spectra. The Laplacian-plus-self-term
interaction family is adopted as \Cref{hyp:gauss-global-interaction}; it does
not follow from a general linear-Gaussian directed model, since an
orientation-reversing gain, and for odd \(K\) a v-structure, leave it in every
chart reached by orientation-preserving frames
(\Cref{prop:gauss-law-level-witnesses}). Within the family, hard aggregation
is closed (\Cref{prop:cg-gaussian-aggregation}) and is the fine law conditioned
on zero within-cluster disagreement (\Cref{prop:cg-aggregation-conditioning}),
although closure selects neither the family nor a nondegenerate partition.
Unrestricted matrix-weighted Kron reduction leaves the family
(\Cref{prop:cg-kron-leaves-family}); the fixed congruence-diagonal cone is
closed (\Cref{thm:cg-congruence-diagonal-kron}), the maximal pair-tested
families at a fixed hidden anchor are the full anchor-congruence-diagonal
cones (\Cref{thm:cg-independent-kron-classification}), and on a fixed flat
graph with independent positive-definite weights and \(K\geq2\), closure for
every weight assignment holds exactly when every hidden component has at most
two retained terminals (\Cref{thm:cg-kron-boundary-classification}).

Coarse families in this supplement admit newly generated positive-semidefinite
self terms. The strict family in which cutting a pair creates no new anchor is
a named subfamily, \Cref{def:cg-no-new-anchor-pair} of the main manuscript,
and the completed-square criterion that separates the two is
\Cref{lem:cg-cut-completed-square}.

Operator coarse-graining is narrower than law-level coarse-graining.
Aggregation is a positive semigroup on the interaction parameters
(\Cref{prop:grg-aggregation-semigroup}), so a rescaling must be declared before
any flow is typed, and on the interior of the cone the rescaled step is a map
of normalized laws (\Cref{prop:grg-interior-law-step}). Bi-additive
parameters are closed under additive pushforward (\Cref{sec:grg-biadditive}),
and under a measurability hypothesis positive additive rules are forced to be
linear (\Cref{prop:grg-positive-additive-linear}). Fixed rays, their attraction, and
independence of the blocking scheme need a hierarchical identification and a
limiting protocol; they belong to Supplement B, and this supplement states no
fixed point of a scale map.

\Cref{ch:expfamily} develops exponential families, their operator layer, and
the information projection. \Cref{ch:generative-gaussian} realizes the
directed law in information form and proves the Schur-complement lemma on
which the later chapters rest. \Cref{ch:gaussian} characterizes when the
belief precision lies in the interaction family, and \Cref{ch:restrictions}
computes the costs of recognition restrictions. \Cref{ch:infogeometry}
and \Cref{ch:gaussian-hermite} treat Fisher geometry, invariants, and the
Hermite spectrum of score tangents. \Cref{ch:gaussian-coarsegraining} and
\Cref{ch:gaussian-nonflat} treat aggregation, Kron closure, nonflat links,
relaxation, and quotient volume; \Cref{ch:gaussian-renormalization} treats the
operator arrow; and \Cref{ch:obstructions} treats reciprocal pairs, anchored
stars, and loop holonomy. \Cref{app:decision-metatheory} holds the decision
metatheory for semialgebraic families and the treatment of singular Gaussian
laws, and \Cref{app:gaussian-backmatter} holds the notation, open problems,
and numerical index of this supplement.

This supplement asserts nothing for a non-Gaussian belief fiber beyond the
exponential-family statements of \Cref{ch:expfamily}, and it states no fixed
point, attraction, or asymptotic scaling result for a scale map.
\status{NOT-CLAIMED}
